\section{Background and Overview}\label{sec:intro_background}
\noindent In recent years, the rapid growth of autonomous driving and connected vehicular networks has accelerated the development of the Internet of Vehicles (IoV).IoV makes it possible for different vehicles to communicate with one another and also with road infrastructures making it as a framework for cooperative intelligent transportation system (C-ITS) \cite{chen2025lbdt}. Applications such as automated collision avoidance, real-time dynamic rerouting (for example, through vehicle-to-vehicle (V2V) communication), cooperative environmental sensing and road hazard detection are based on the constant sharing of sensor and telemetry data \cite{javaid2020scalable}.

However, data sharing in vehicular networks cannot be based only on goodwill. Data gathering uses local computational and storage resources as well as energy in the on-board units (OBUs). That is why cars in vehicular networks act in a self-interested way and need incentives for sharing data \cite{kang2019toward}. Data trading protocols enable data-buyers (e.g.. autonomous vehicles needing real-time road condition updates) to buy verified data from data-sellers (e.g., cars/vehiceles equipped with LiDAR or camera arrays).

Conventional vehicle data markets depend on centralized hardware that connects sellers, buyers, and allows transactions. While it is easy to implement, having a centralized server raises the possibility of failure, increases communication delays across hops in wireless networks, exposes the whereabouts of cars, and creates chances of platform manipulation \cite{michelin2018speedychain}. Existing literature advocates for decentralized trading based on blockchain technology as a new solution that offers reliable and transparent waiting for payments without the use of any centralized parties \cite{liu2018computation, firdaus2021blockchain}.

\section{Motivation and Key Challenges}\label{sec:intro_motivation}
\noindent Deploying standard blockchain technologies directly in high-mobility IoV environments introduces several severe technical bottlenecks:

\begin{enumerate}
    \item \textbf{Transaction Confirmation Latency vs. Mobility Dynamics:} Standard Proof-of-Work (PoW) or traditional Proof-of-Stake (PoS) blockchains suffer from low transaction throughput (TPS) and multi-minute confirmation delays \cite{nakamoto2008bitcoin}. In vehicular networks, where topologies change rapidly within seconds, transaction latencies exceeding a few seconds cause market mismatches and outdated data delivery.
    \item \textbf{Dishonest Behaviors and Identity Spoofing:} Malicious or opportunistic vehicular nodes may transmit fabricated, noisy, or incomplete data to maximize profits. Without an adaptive trust evaluation model, bad actors can exploit trading protocols through selective packet dropping, free-riding, or collusion attacks \cite{rosenstatter2018modelling}.
    \item \textbf{Market Inefficiency and Unfair Pricing:} Traditional auction protocols often fail to account for data freshness (latency) and seller reputation simultaneously. Unregulated pricing can discourage honest participants while rewarding dishonest sellers who offer low-quality data at predatory prices \cite{guan2020towards}.
    \item \textbf{High Consensus Communication Overhead:} Classical Byzantine Fault Tolerance (BFT) consensus protocols require $O(N^2)$ message complexity during commit rounds \cite{lamport1982byzantine}. In dense urban intersections with hundreds of participating Roadside Units (RSUs), this communication burden saturates the wireless channel and causes severe network congestion.
\end{enumerate}

\section{Problem Statement}\label{sec:intro_problem}
\noindent With the dynamic nature of the vehicular network that features time-dependent topology, changing vehicle density, and maybe compromised nodes, this project will focus on the creation, achievement, and assessment of the lightweight blockchain-based data trading method (LBDT). Specifically, the method must ensure a low confirmation time, high transaction capacity, solid protection against malicious actions, just market value, and scalable consensus under strict vehicle mobility limitations.

\section{Project Objectives}\label{sec:intro_objectives}
\noindent In order to address the issues described above this project has the following specific technical goals:
\begin{enumerate}
    \item \textbf{Dual Parallel Chain Architecture:} Create a separate architecture for Keyblocks (keeping epoch setups and the global reputation list) and Microblocks (sensing localized high productivity data trades).
    \item \textbf{Aged Gompertz Reputation \& ML Classifier:} Create a realtime model which is based on the transaction frequency, aging effects, and a ML-based Logistic Regression classifier ($p_{honest}$) to estimate and weed out bad performing nodes.
    \item \textbf{Reputation-Aware Double Auction:} Create a double auction protocol that would adjust the selling price of sellers based on the execution time and node reputation in order to provide honest market clearing prices.
    \item \textbf{Lightweight EC-VRF \& Gosig BFT Consensus:} Use Elliptic Curve Verifiable Random Function (EC-VRF) for unpredictable committee leader selection plus high performance Gosig BFT consensus protocol which limits the communication overhead to ($5N-3$) messages.
    \item \textbf{Dynamic Zone Sharding:} Use dynamic geographic zone rules based on utility load scores ($U_z$) to shift overutilized clusters and merge underutilized zones.
    \item \textbf{Co-Simulation Framework:} Ensure the functionality of the developed frame-
    work by combining the Python logic with C++ OMNeT++/Veins simulation and simulation program SUMO.
\end{enumerate}

\section{Major Technical Contributions}\label{sec:intro_contributions}
\noindent This report summarizes the major technical contributions of the project, including the following developments:
\begin{itemize}
    \item An Aged Gompertz Reputation Model which has been developed that has enabled log transformed interaction counts and an adjustable parameter adjustment method with the Logistic Regression classifier achieving the empirical ROC-AUC of 0.9991 for 10,400 dataset samples.
    \item A Double Auction matching algorithm that has been developed including a penalty function for bid losses such as  $loss = \alpha_1 (\Delta t)^{\beta_1} + \alpha_2 (1 - Rep_i)^{\beta_2}$, ensuring that sellers characterized by low reputation or by high latency are penalized appropriately.
    \item A consensus pipeline has been created that includes the SECP256k1 EC-VRF leader election as well as the four-phase Gosig BFT consensus engine limiting the use of messages to ($5N – 3$) per step of consensus.
    \item A thorough co-simulation pipeline has been created that allows for SUMO, Veins, OMNeT++, and Python to be implemented, which demonstrates an average transaction confirmation time of 2.4828 seconds and limit on the malicious reputation levels of the committee below 19.4\% (well below the permissible level of 33.3\% with respect to BFT).




\end{itemize}

\section{Organization of the Report}\label{sec:intro_organization}

\noindent The remainder of this report is organized as follows:

{\emergencystretch=2em
\begin{itemize}
    \item \textbf{Chapter 2 (Literature Review):} Reviews prior art in vehicular networks, blockchain-based data trading, reputation systems, double auctions, and BFT consensus.

    \item \textbf{Chapter 3 (Problem Statement \& System Model):} Formulates the formal mathematical model, system components, and security threat model.

    \item \textbf{Chapter 4 (Implemented LBDT Scheme \& Methodology):} Presents the technical design and mathematical formulation of the dual parallel chain, reputation model, ML classifier, double auction engine, EC-VRF Gosig BFT, and dynamic sharding logic.

    \item \textbf{Chapter 5 (Experiment and Results):} Details the experimental setup, simulation environment, empirical performance metrics, and comparative evaluation graphs.

    \item \textbf{Chapter 6 (Conclusions and Future Scope):} Concludes the report with a summary of findings, limitations, and future research directions.
\end{itemize}
}